% titel: Blatt 8, A1: Python-Code
% skript: Kap. 4.3 (Newton-Verfahren), Kap. 5 (Summen)

% Teil I: Werkzeuge, kein Lösungsweg.

Zwei kleine Python-Schnipsel verstehen und mathematisch übersetzen: eine
Summe über Potenzen sowie eine Näherungsschleife, die eine Kubikwurzel
berechnet.

\begin{werkzeugkasten}[Summe als mathematischer Ausdruck]
Der Ausdruck
\begin{verbatim}
sum([1/(k**p) for k in range(1, n+1)])
\end{verbatim}
summiert über $k=1,\dots,n$ jeweils den Term $1/k^p$:
\begin{equation*}
  \sum_{k=1}^{n}\frac{1}{k^p}.
\end{equation*}
Für $p=1$ handelt es sich um die (partielle) harmonische Reihe, die langsam
gegen $\infty$ wächst; für $p>1$ konvergiert die Reihe (Zeta-Funktion).

\textit{Skript:} Kap. 5 (Summen).
\end{werkzeugkasten}

\begin{werkzeugkasten}[Newton-Schleife zum Kubikwurzel-Ziehen]
\begin{verbatim}
x = 3
for n in range(11):
    print(x)
    x = x - (x**3 - 3) / (3*x**2)
\end{verbatim}
Der Code startet bei $x=3$ und wendet $11$-mal die Newton-Iteration für
$f(x)=x^3-3$ mit $f'(x)=3x^2$ an:
\begin{equation*}
  x_{k+1}=x_k-\frac{x_k^3-3}{3x_k^2}.
\end{equation*}
Dabei ist $x_k^3-3$ der Fehler zu $f(x)=x^3-3=0$, gesucht ist die reelle
Lösung $x=\sqrt[3]{3}\approx1{,}442$.

\textit{Skript:} Kap. 4.3 (Newton-Verfahren).
\end{werkzeugkasten}

\begin{achtung}[Darauf achten]
\begin{itemize}
  \item \texttt{range(1,n+1)} läuft von $1$ bis $n$ (den Endwert
        $n$ inklusive), \texttt{range(n)} von $0$ bis $n-1$ -- zählen
        unterscheidet die beiden.
  \item \texttt{k**p} ist Potenz, nicht $k$ mal $p$.
  \item Die Newton-Formel muss mit der Ableitung
        $f'(x)=3x^2$ übereinstimmen; sie ist korrekt, wenn
        $\frac{x_k^3-3}{3x_k^2}$ gegen $0$ strebt (dann konvergiert $x_k$).
  \item \texttt{print(x)} vor der Aktualisierung druckt das $x$, \emph{bevor}
        es verändert wird.
\end{itemize}
\end{achtung}
